\section{Relación entre las capacidades fundamentales y el post-entrenamiento}

\subsection{El paradigma de dos etapas}

Agentic Vision revela un paradigma de entrenamiento importante:

\begin{importantbox}{Capacidades fundamentales + liberación mediante post-entrenamiento}
\begin{enumerate}
    \item \textbf{Etapa de preentrenamiento}: se entrenan las \textbf{capacidades visuales fundamentales} del modelo: detección, segmentación, reconocimiento. Estas capacidades están «dormidas»; el modelo ya sabe cómo hacerlo, pero no sabe cuándo hacerlo ni cómo combinarlo.
    \item \textbf{Etapa de post-entrenamiento}: se encapsula una herramienta (como Code Execution) y, mediante aprendizaje por refuerzo, se incentiva al modelo a formar un patrón de comportamiento de \textbf{llamadas a herramientas en múltiples turnos}, con lo que se «libera» la capacidad fundamental.
\end{enumerate}
\end{importantbox}

Este paradigma ofrece dos enseñanzas para la investigación de agentes de IA:

\begin{enumerate}
    \item \textbf{Para quienes diseñan prompts/agentes}: mediante un diseño ingenioso de Visual Prompting y de herramientas, es posible que tareas antes inviables se vuelvan viables sin cambiar los parámetros del modelo
    \item \textbf{Para quienes entrenan modelos}: la calidad de las capacidades fundamentales determina el techo del agente después del post-entrenamiento
\end{enumerate}

\subsection{El papel del Trigger Prompt}

En la demostración oficial de Gemini, el prompt incluye ciertas «palabras disparadoras» para activar la capacidad de Agentic Vision:

\begin{itemize}
    \item "crop out all ..."
    \item "zoom in to see ..."
    \item "annotate on the image ..."
\end{itemize}

\begin{knowledgebox}{El Trigger Prompt es un diseño transitorio}
Si el entrenamiento del modelo fuera suficientemente bueno, en teoría no harían falta estos Trigger Prompt: el modelo debería poder decidir por sí mismo cuándo necesita hacer zoom in y cuándo necesita anotar. La existencia del Trigger Prompt indica que el post-entrenamiento del modelo actual todavía tiene margen de optimización.
\end{knowledgebox}

\subsection{Resumen de la sección}

La enseñanza central de Agentic Vision es que el techo de capacidad de un agente lo determinan las capacidades fundamentales del modelo, y la función del post-entrenamiento es combinar y liberar esas capacidades fundamentales como un comportamiento de razonamiento completo en múltiples turnos.

